\chapter{Admission существующего носителя charged mediator}
\label{ch:v8-charged-mediator-existing-carrier-admission}

\section{Уточнение семейного разложения}

Предыдущий coarse-аудит применил стандартную семейную тройку ко всему
трёхкратному charged-singlet типу и получил нулевой fixed layer. Однако
ранние endpoint-гейты уже различили grading-компоненты:
\begin{equation}
 H_{e_R}^{\rm old}=\mathbb C e_R^{(s)}
 \oplus\operatorname{span}_{\mathbb C}
 \{e_R^{(t,1)},e_R^{(t,2)}\},
 \qquad \Gamma=(-1)\oplus I_2.
 \label{eq:v8-mediator-existing-old-decomposition}
\end{equation}
Линия $e_R^{(s)}$ является семейным синглетом. Две положительно
градуированные линии являются не полной тройкой, а её незавершённой частью.
Условное минимальное глобальное замыкание добавляет $e_R^{(t,0)}$:
\begin{equation}
 H_{e_R}^{\rm ext}=\mathbb C e_R^{(s)}\oplus\mathbb C^3_t
 \simeq\mathbf1\oplus\mathbf3,
 \qquad \dim H_{e_R}^{\rm ext}=4.
 \label{eq:v8-mediator-existing-one-plus-three}
\end{equation}
Для семейных генераторов
\begin{equation}
 J_a^{\rm ext}=0\oplus J_a^{(\mathbf3)}
 \label{eq:v8-mediator-existing-family-action}
\end{equation}
общий fixed layer одномерен:
\begin{equation}
 \bigcap_{a=1}^{3}\ker J_a^{\rm ext}
 =\mathbb C e_R^{(s)}.
 \label{eq:v8-mediator-existing-fixed-line}
\end{equation}
Тем самым прежнее утверждение «все три старые линии образуют irreducible
triplet» пересмотрено. Нулевая fixed-line относилась к неверно укрупнённому
представлению.

\section{Grading выбирает единственный mediator}

Нейтральное расширение содержит
\begin{equation}
 H_{\rm n}=\mathbb C s_0\oplus\mathbb C a_0,
 \qquad \Gamma_{\rm n}=\operatorname{diag}(+1,-1),
 \qquad Y_{\rm n}=0.
 \label{eq:v8-mediator-existing-neutral-source}
\end{equation}
Ищем family-invariant нечётную карту $T:H_{\rm n}\to H_{e_R}^{\rm ext}$:
\begin{equation}
 J_a^{\rm ext}T=0,\qquad
 \Gamma_{e_R}^{\rm ext}T+T\Gamma_{\rm n}=0.
 \label{eq:v8-mediator-existing-intertwiner-system}
\end{equation}
Точная линейная система имеет восемь неизвестных, ранг семь и потому
\begin{equation}
 \dim_{\mathbb C}\operatorname{Hom}_{SO(3)}^{\rm odd}
 (H_{\rm n},H_{e_R}^{\rm ext})=1,
 \qquad
 T=\lambda |e_R^{(s)}\rangle\langle s_0|.
 \label{eq:v8-mediator-existing-unique-odd-hom}
\end{equation}
Поскольку $Y e_R^{(s)}=-e_R^{(s)}$, получаем
\begin{equation}
 [Y,T]=-T.
 \label{eq:v8-mediator-existing-charge-minus}
\end{equation}
Real-сопряжённая карта $T^c=JTJ^{-1}$ имеет заряд $+1$ и предоставляет
нужный environment-transition $b_+$. Следовательно, независимую новую
charged-singlet пару добавлять не требуется: она уже присутствует в
Real-удвоении старой singlet-линии.

\section{Совпадение с условным 47-кадром}

Карта $T$ является ровно ветвью $z_0$ ранее классифицированного семейства
\begin{equation}
 V(z)=z_0|e_R^{(s)}\rangle\langle s_0|
 +z_1|e_R^{(t,1)}\rangle\langle a_0|
 +z_2|e_R^{(t,2)}\rangle\langle a_0|.
 \label{eq:v8-mediator-existing-vz-family}
\end{equation}
Family-invariance и grading совместно сокращают $\mathbb{RP}^2$ до
единственной точки $[1:0:0]$. Две самосопряжённые квадратуры
\begin{equation}
 C_1=T+T^\dagger,\qquad
 C_2=i(T-T^\dagger)
 \label{eq:v8-mediator-existing-quadratures}
\end{equation}
образуют минимальную $U(1)$-орбиту:
\begin{equation}
 i[Y,C_1]=-C_2,\qquad i[Y,C_2]=C_1,
 \qquad
 \bigl(\Tr C_r^\dagger C_s\bigr)_{r,s}=2I_2.
 \label{eq:v8-mediator-existing-u1-gram}
\end{equation}
Именно эти направления входят в условный $47$-кадр. Но $45$-кадр их ещё
не содержит.

\section{Уровни admission}

На сыром $H_{21}$ имеются charged family-singlet линия и её Real-партнёр,
но нет нейтрального $s_0$. Поэтому
\begin{equation}
 N_{\rm raw\ H21}=2/7.
 \label{eq:v8-mediator-existing-raw-ledger}
\end{equation}
После условных расширений $H_{23}$, $H_{24}$ и frame $45\to47$ все
структурные требования к carrier выполнены:
\begin{equation}
 N_{\rm conditional\ carrier}=7/7.
 \label{eq:v8-mediator-existing-conditional-ledger}
\end{equation}
Однако сами расширения не происходят из текущего parent action. Кроме того,
не выведены резонанс, состояние среды, coupling и механизм обновления:
\begin{equation}
 N_{\rm dynamic\ origin}=0/5:
 \quad\{7\gamma,\rho_E=|0\rangle\langle0|,g,
 \text{fresh chain},\text{clock--rate anchor}\}.
 \label{eq:v8-mediator-existing-dynamic-ledger}
\end{equation}

\begin{theorem}[Условный admission без нового charged типа]
После точного разложения charged multiplicity как
$\mathbf1\oplus\mathbf3$ family-invariant нечётный Hom из нейтральной пары
одномерен и порождён $|e_R^{(s)}\rangle\langle s_0|$. Его Real-сопряжение
имеет требуемый заряд $+1$. Поэтому условные расширенные carriers допускают
минимальный mediator без ещё одной charged-singlet пары.
\end{theorem}

\begin{nogocriterion}[Admission не есть origin]
Наличие старой singlet-линии не создаёт нейтральный вакуум $s_0$, глобальное
семейное замыкание, две квадратуры $47$-кадра, резонансную энергию или
слабый collision-scale. Нельзя переносить условный carrier-pass на текущий
parent action: его raw admission остаётся $2/7$.
\end{nogocriterion}

\begin{protocol}\begin{enumerate}
\item старых charged-singlet линий: \textbf{$3$};
\item все они образуют irreducible triplet: \textbf{нет};
\item точное условное разложение: \textbf{$\mathbf1\oplus\mathbf3$};
\item family-fixed charged dimension: \textbf{$1$};
\item единственный нечётный invariant mediator: \textbf{$|e_R^{(s)}\rangle\langle s_0|$};
\item его заряд: \textbf{$-1$};
\item заряд Real-сопряжённой карты: \textbf{$+1$};
\item минимальных Hermitian quadratures: \textbf{$2$};
\item raw $H_{21}$ admission: \textbf{$2/7$};
\item conditional extended admission: \textbf{$7/7$};
\item dynamic parent-origin: \textbf{$0/5$};
\item следующий гейт: \textbf{parent-origin энергии, состояния и скорости}.
\end{enumerate}\end{protocol}

Машинный сертификат сохранён в
\path{s2t_v8_baryon_c0_singlet_triplet_central_gap_isotypic_channel_extra_edge_mass_two_source_cubic_incidence_up_down_binary_selector_minimal_typed_discriminator_binary_carrier_charged_environment_mediator_existing_carrier_admission_gate_results.json}.

\begin{thebibliography}{9}
\bibitem{v8-mediator-existing-stinespring} W.~F.~Stinespring,
\emph{Positive Functions on C*-Algebras},
Proc. Amer. Math. Soc. 6 (1955) 211.
\bibitem{v8-mediator-existing-haapasalo} E.~Haapasalo,
\emph{The Choi--Jamio\l{}kowski Isomorphism and Covariant Quantum Channels},
Quantum Stud. Math. Found. 8 (2021) 1; arXiv:1906.11442.
\end{thebibliography}